\documentclass[a4paper,12pt]{ctexart}
\usepackage{xcolor}
\usepackage[top=1.8cm, bottom=1.8cm, left=1.8cm, right=1.8cm]{geometry}
\usepackage{array}
\usepackage{longtable}
\usepackage{hyperref}
\hypersetup{colorlinks=true,linkcolor=blue!70!black}
\linespread{1.35}

\definecolor{donegreen}{RGB}{40,140,80}
\definecolor{economics}{RGB}{180,60,30}
\definecolor{philosophy}{RGB}{30,100,160}
\definecolor{politics}{RGB}{20,130,60}
\definecolor{law}{RGB}{140,90,40}
\definecolor{ethics}{RGB}{120,50,120}

\newcommand{\tageconomics}{\textcolor{economics}{\small [经济]}}
\newcommand{\tagphilosophy}{\textcolor{philosophy}{\small [哲学]}}
\newcommand{\tagpolitics}{\textcolor{politics}{\small [政治]}}
\newcommand{\taglaw}{\textcolor{law}{\small [法律]}}
\newcommand{\tagethics}{\textcolor{ethics}{\small [伦理]}}
\newcommand{\done}{\textcolor{donegreen}{\textbf{【已写】}}}

\title{\textcolor{blue!70!black}{\Huge\textbf{社会科学小故事 · 课程大纲}}}
\author{}
\date{}

\begin{document}
\maketitle
\thispagestyle{empty}

\section*{\textcolor{blue!70!black}{设计思路}}

本模块面向北京地区四年级起点的孩子，用讲故事的口吻介绍社会科学。
每课围绕一个孩子生活里能接触到、能追问、能讨论的社会现象展开，
从具体体验过渡到背后的社会科学道理。不做概念背诵，不要求记忆，
只求唤起对社会运转逻辑的好奇心。

选课的三个原则：
\begin{enumerate}
  \item \textbf{从身边出发}——孩子见过、经历过、听说过
  \item \textbf{可以追问"为什么"}——每个现象背后有一条清晰的分析线索
  \item \textbf{学科均衡}——经济、政治、法律、哲学、伦理各有覆盖
\end{enumerate}

\vspace{0.3cm}
\textcolor{gray}{\small 当前规划20课，可根据实际教学进度增删。}

\newpage

\section*{\textcolor{orange!80!black}{一、经济（钱与交换的故事）}}

\begin{enumerate}

\item[\textbf{第1课}] \textbf{铅笔的故事——看不见的手} \done \\
\tageconomics \\
\textbf{内容：} 一支最普通的铅笔——木头、石墨、橡皮、金属箍、油漆——
没有任何一个人能独立制造出它。成千上万的人（林场工人、石墨矿工、
橡胶种植园主、化工厂工人、运输司机）在全球各地合作，只为了让你
花几毛钱买到一支笔。没有人下令"做铅笔"，但市场把这所有的事情
都安排好了。引出"自发秩序"和"看不见的手"的概念。\\
\textbf{推荐理由：} 铅笔是孩子每天用的东西，但几乎没人想过"它背后的
全球合作网"。这是经济学最经典的开场白。

\item[\textbf{第2课}] \textbf{为什么牛奶装在方盒子里而可乐装在圆瓶子里？} \\
\tageconomics \\
\textbf{内容：} 超市的货架上——牛奶是方盒子、可乐是圆瓶子、薯片是
鼓鼓的袋子。为什么？因为牛奶需要冷藏，方盒子节省冷柜空间；
可乐在常温货架上，圆瓶拿起来更舒服；薯片怕碎，鼓袋子装氮气
防震。引出"成本与收益的平衡"——每一个包装都在省钱和好用之间
找平衡。\\
\textbf{推荐理由：} 孩子逛超市时都能看到，从"为什么长这样"追问
经济逻辑。

\item[\textbf{第3课}] \textbf{零花钱怎么花？——选择就是放弃} \\
\tageconomics \\
\textbf{内容：} 你每周有20块零花钱。买玩具就不能买零食，买贴纸
就不能买漫画。"选择一样意味着放弃另一样"——这就是"机会成本"。
扩展到更大的决策：国家修高铁不修公园、你花时间玩游戏而不学
英语——每种选择的背后都藏着被放弃的选项。\\
\textbf{推荐理由：} 每个孩子都有支配零花钱的体验，最自然的
经济学入口。

\item[\textbf{第4课}] \textbf{以物易物的麻烦——货币为什么诞生} \\
\tageconomics \\
\textbf{内容：} 原始社会里，你想用自己家的羊换邻居家的斧头。
但邻居不想要羊，他想要米。你得先找到想要羊的人换到米，
再用米换斧头。太麻烦了！于是人们想出一个办法：找一个"人人
都想要"的东西来当中间人——贝壳、盐、金银、最后是纸币和
手机里的数字。引出货币的三大功能：价值尺度、交换媒介、
储藏手段。\\
\textbf{推荐理由：} 拿东西换东西是孩子幼儿园时代的"交易"模式，
从他们的亲身经历推进到货币本质。

\item[\textbf{第5课}] \textbf{为什么可乐涨价了？——供需跷跷板} \\
\tageconomics \\
\textbf{内容：} 去年可乐3块，今年变成3块5。为什么？夏天热，
喝可乐的人多（需求↑），但工厂来不及生产那么多（供给不变），
价格就涨了。反过来，如果大家都改喝茶了（需求↓），可乐卖不
出去，超市就会打折。引出"供给和需求"以及它们如何决定价格。
延伸到演唱会门票为什么那么贵、口罩为什么在疫情期间涨价。\\
\textbf{推荐理由：} 孩子每次买东西都能观察到的价格变化，背后
是经济学最核心的模型。

\end{enumerate}

\section*{\textcolor{orange!80!black}{二、政治（我们怎么一起生活）}}

\begin{enumerate}

\item[\textbf{第6课}] \textbf{为什么要有红绿灯？——规则的必要性} \\
\tagpolitics \\
\textbf{内容：} 十字路口没有红绿灯会怎样？——大家抢着走、
谁也走不了、撞车。有了红绿灯，虽然有时要等，但所有人都能
更安全更快地通过。引出"规则不是为了限制自由，而是为了
保障更大的自由"。延伸到为什么要有法律、学校为什么要定
校规。\\
\textbf{推荐理由：} 每个孩子都等过红绿灯，从最直观的规则体验
开始理解"社会契约"的雏形。

\item[\textbf{第7课}] \textbf{三个和尚没水喝——集体行动难题} \\
\tagpolitics \\
\textbf{内容：} 一个和尚挑水喝，两个和尚抬水喝，三个和尚
没水喝——人多了为什么反而没水喝？引出"搭便车"问题：
每个人都指望别人去挑水，自己坐享其成。延伸到更广的例子：
全班大扫除时有人偷懒、公共区域的灯没人关、税收为什么
不能自愿缴纳。讨论可能的解决方式：分工、监督、激励。\\
\textbf{推荐理由：} 经典故事每个孩子都听过，但没想过它背后
是政治学和政治经济学的核心问题。

\item[\textbf{第8课}] \textbf{投票——少数服从多数公平吗？} \\
\tagpolitics \\
\textbf{内容：} 班里决定周末去动物园还是博物馆。投票：12人
去动物园，11人去博物馆。动物园赢了。但11个人不开心——
"多数人"的决定对"少数人"公平吗？引出投票制度的本质：
它不是"找到正确答案"的方法，而是"大家都接受"的方法。
延伸到选举制度、为什么有些国家用"排名投票"替代"简单
多数"。\\
\textbf{推荐理由：} 孩子在班级里经常投票，这是让他们思考
"民主的难题"的最佳起点。

\end{enumerate}

\section*{\textcolor{orange!80!black}{三、法律（对与错的边界）}}

\begin{enumerate}

\item[\textbf{第9课}] \textbf{分蛋糕——切的人最后选} \\
\taglaw \\
\textbf{内容：} 两个小朋友分一块蛋糕。怎么分最公平？——
一个人切，另一个人先选。切的人为了自己不拿到最小的那块，
会尽量切得一样大。引出"程序正义"的概念：结果公平靠的不是
切蛋糕的人品，而是分蛋糕的规则设计。延伸到法律中的"回避
制度"（法官不能审自己家人的案子）、考试中的"双盲评分"。\\
\textbf{推荐理由：} 每个孩子都分过东西，"切的人最后选"是
他们能直观理解的制度设计智慧。

\item[\textbf{第10课}] \textbf{偷一块钱和偷一万块钱有什么区别？} \\
\taglaw \\
\textbf{内容：} 同样是偷东西，为什么偷一万块判得更重？
引出法律的几个原则：(1) 罪刑相适应——刑罚的轻重和违法
的严重程度匹配；(2) 故意和过失的区别——不小心踩了别人的
脚和故意踩一脚，处理方式完全不同；(3) 年龄——为什么法律
对小孩和成年人的要求不一样。用孩子能理解的例子讲这三个
概念。\\
\textbf{推荐理由：} 孩子天然有"公平感"，从"一样的事情为
什么处罚不一样"的问题出发，理解法律的精细逻辑。

\item[\textbf{第11课}] \textbf{谁的错？——一起简单的"交通事故"} \\
\taglaw \\
\textbf{内容：} 小明在操场上跑，撞倒了一个同学，同学的膝盖
破了。谁的责任？——如果是小明故意推的，是故意的错；
如果小明在追球没看路，是过失；如果同学突然从墙角冲出来，
可能双方都有责任。引出"因果关系"和"责任分配"的法律
思维。延伸到更复杂的例子：一个司机正常行驶，突然有人闯
红灯被撞——法律怎么判？\\
\textbf{推荐理由：} 孩子在学校天天发生"意外"，用他们的
日常冲突引出法律中的因果推理。

\end{enumerate}

\section*{\textcolor{orange!80!black}{四、哲学（问"为什么"的学问）}}

\begin{enumerate}

\item[\textbf{第12课}] \textbf{苏格拉底的问题——为什么要问"为什么"} \\
\tagphilosophy \\
\textbf{内容：} 苏格拉底在雅典街头不停地问别人问题："什么是
勇敢？""什么是正义？"他问得那些人哑口无言。他不是为了
为难别人——他是认为"不经过审视的人生不值得过"。引出
"批判性思维"的起点：对任何理所当然的说法都问一句"你
凭什么这么说？"。\\
\textbf{推荐理由：} 每个孩子都是天生的"为什么机器"，苏格拉底
是他们精神上的同类。把爱追问从惹人烦变成受尊重。

\item[\textbf{第13课}] \textbf{忒修斯之船——你还是原来的你吗？} \\
\tagphilosophy \\
\textbf{内容：} 一艘船在海上航行了几百年，木板一块块烂掉
被替换。等到所有木板都被换了一遍——这艘船还是原来的那艘
船吗？如果把你身体的细胞全部替换一遍——你还是你吗？
你觉得"你"指的是什么——你的身体？你的记忆？你的性格？
引出"同一性"这个哲学问题。\\
\textbf{推荐理由：} 问题本身就像一个谜语，孩子会觉得好玩，
然后发现越想越深——这就是哲学的魔力。

\item[\textbf{第14课}] \textbf{自由是什么？——不做作业的自由和过马路的自由} \\
\tagphilosophy \\
\textbf{内容：} "自由"是不是想做什么就做什么？如果是，
那小明想打小张，他的"自由"和小张"不被挨打的自由"冲突了。
引出约翰·密尔的"伤害原则"：你的自由止于别人的鼻子。
延伸到：为什么不系安全带要被罚款（你觉得的"自己的事"
可能对社会造成负担）、为什么言论自由不意味着可以造谣。\\
\textbf{推荐理由：} 孩子对"自由"有强烈的直觉——"凭什么
管我"。引导他们思考自由的边界在哪里。

\end{enumerate}

\section*{\textcolor{orange!80!black}{五、伦理（什么是对、什么是错）}}

\begin{enumerate}

\item[\textbf{第15课}] \textbf{电车难题——没有完美答案的问题} \\
\tagethics \\
\textbf{内容：} 一辆失控的电车冲向五个被绑在轨道上的人。
你可以拉杆让电车转到另一条轨道上——但那条轨道上有一个
人。你拉不拉？没有一个"正确答案"，但不同的回答展示了
不同的道德原则：(1) 功利主义——救五个牺牲一个，总损失
最小；(2) 道义论——故意杀死一个人是错的，无论结果如何；
(3) 你希望别人怎么选？\\
\textbf{推荐理由：} 孩子在做道德判断时本能地快速下结论，
电车难题迫使他们慢下来思考"做决定的理由是什么"。

\item[\textbf{第16课}] \textbf{不公平的起跑线——运气和努力哪个更重要？} \\
\tagethics \\
\textbf{内容：} 两个孩子一起跑步比赛：一个穿着跑鞋在塑胶跑道上，
一个光着脚在沙地上。公平吗？但生活中的"起跑线"差异更大：
有的孩子出生在大城市、父母是高知；有的在偏远农村、父母
常年在外打工——这不是他们自己选的。引出罗尔斯的"无知之幕"：
如果你不知道自己会出生在什么样的家庭，你希望社会是什么样的？\\
\textbf{推荐理由：} 孩子对"不公平"非常敏感。从他们能理解的
"起跑线差异"出发，引出对社会正义的思考。

\item[\textbf{第17课}] \textbf{要不要说实话？——诚实与善意的谎言} \\
\tagethics \\
\textbf{内容：} 奶奶织了一件毛衣给你，说实话你不喜欢——你该
说"好难看"吗？朋友新剪了头发问你怎么样，你觉得不好看——
怎么说？引出"道德冲突"：诚实是一个好原则，但有时候为了
不伤害别人的感受，人们选择"善意的谎言"。这算撒谎吗？
什么时候可以说谎？有没有永远不能说谎的原则？\\
\textbf{推荐理由：} 孩子每天都在面对"要不要说实话"的场景。
不是教他们标准答案，而是教他们理解道德决定中的权衡。

\end{enumerate}

\newpage

\section*{\textcolor{blue!70!black}{学科分布统计}}

\begin{center}
\begin{tabular}{|l|c|c|c|}
\hline
\textbf{学科方向} & \textbf{课号} & \textbf{课数} & \textbf{占比} \\
\hline
经济 & 第1-5课 & 5 & 29\% \\
政治 & 第6-8课 & 3 & 18\% \\
法律 & 第9-11课 & 3 & 18\% \\
哲学 & 第12-14课 & 3 & 18\% \\
伦理 & 第15-17课 & 3 & 18\% \\
\hline
\textbf{合计} & & \textbf{17课} & \\
\hline
\end{tabular}
\end{center}

\vspace{0.3cm}

\textcolor{gray}{\small 注：第1课（铅笔的故事）已写，其余待写。第15课电车难题兼具哲学与伦理特征，归入伦理类。}

\vspace{0.5cm}

\begin{center}
\rule{0.5\textwidth}{0.5pt}
\vspace{0.3cm}
\textcolor{gray}{\small 大纲版本：2026-05-25 · 初稿}
\end{center}

\end{document}
